\section{Introduction}
\label{sec:intro}
% \vspace{-10px}
A long-standing goal in robotics is to build generalist manipulation policies that can follow natural-language instructions and manipulate arbitrary objects across diverse scenes~\citep{kim2024openvla, openvla_oft, kim2026cosmos, black2024pi_0}. To achieve this, a general manipulation model must not only recognize objects and parse instructions, but also reason about how the physical world will evolve under its own actions. This requires a unified understanding of language, visual appearance, scene geometry, robot state, and physical dynamics. 

Recent progress has therefore increasingly relied on large-scale foundation models as pretrained substrates for robot policies. Vision-language-action models (VLAs) build on vision-language models whose representations are aligned with natural language, and learn to map visual and linguistic tokens to robot actions~\cite{zitkovich2023rt,kim2024openvla,openvla_oft,black2024pi_0,pi05,pi0_fast}. Video world-action models (WAMs) instead leverage pretrained video generation models, using their world prediction priors to jointly model future frames and actions~\cite{kim2026cosmos,pai2025mimic}. While these approaches have shown impressive language-conditioned manipulation ability, they are fundamentally in 2D: 3D cues such as depth, scale, and occlusion are left implicit in monocular cues that the action decoder must disentangle on its own, leading to limited generalization across environment changes, especially in robot initial state and camera viewpoint~\cite{fei2025libero, wilcox2025adapt3r}. 

To overcome this limitation, recent work incorporates 3D geometric information into robot policies. One line learns policies directly on explicit 3D observations such as raw point clouds~\citep{ze20243d,huang2026pointworld}, demonstrating the value of geometry for generalization but typically requiring task-specific encoders trained from scratch. With the emergence of Geometric Foundation Models (GFMs)~\cite{lin2025depth,wang2025vggt,wang2026vggt}, some works transfer pretrained geometric priors into VLA policies, either distilling selected GFM features into the VLA backbone through representation alignment~\citep{yu2024representation,li2025spatial,sun2026rocket} or attaching a lightweight action head on top of a GFM's final features~\citep{qian2025gp3,ge2025vggt}. These improve spatial awareness, but use the GFM only as a static feature extractor: its multi-layer geometric structure is never repurposed as the policy's own temporal and action-generating substrate.  

In this work, we propose \textbf{Geometric Action Model (GAM)}, which directly repurposes a GFM as a manipulation policy by using it as a shared medium for perception, future-state prediction, and action decoding. We show that by jointly predicting future action and geometry, geometric world dynamics can be inherently incorporated into robot policies.

Specifically, we split the pretrained GFM at an intermediate layer: the shallow layers serve as an observation encoder, while the remaining layers serve as a decoder block. Given the current visual observation, the observation encoder extracts spatially meaningful scene representations. To model how the world evolves over time, we insert a causal transformer at the intermediate layer that predicts future feature representations. This predictor is conditioned on task language, proprioception, and action history by introducing them as additional tokens at each timestep. The predicted future-state tokens are then processed by the remaining GFM decoder together with an action token, allowing the backbone to produce both future geometry and robot actions. An intuitive comparison between existing paradigms and our proposed framework is illustrated in Figure~\ref{fig:intuition}.


Across diverse simulation~\citep{liu2023libero,fei2025libero,nasiriany2024robocasa} and real-world benchmarks, GAM matches or exceeds the success rate of current foundation-model-scale baselines such as VLAs and WAMs while using substantially fewer trainable parameters and substantially faster inference \textbf{(55$\times$ faster)}, and shows improved generalization to unseen scenarios. GAM especially achieves outstanding performance in camera perturbation settings \textbf{($\uparrow$9.7\%p)}, which requires geometric understanding priors. 

Our contributions are as follows:
\vspace{-2.5pt}
\begin{itemize}
    \item We introduce \textbf{Geometric Action Model (GAM)}, a manipulation policy that combines temporal world modeling, latent feature-space prediction, and a geometric foundation-model substrate in a single shared-backbone architecture.
    \item We show that action and geometry can be predicted in a shared token space: a single autoregressive sequence and a single backbone forward pass produce both action tokens and future-scene tokens, decoded by a lightweight action regression head and a depth head.
    \item We demonstrate that across diverse simulation and real-world manipulation benchmarks, GAM is simultaneously more accurate, more robust, faster, and lighter than current foundation-model-scale alternatives.
\end{itemize}